% Undecidable quantum storage: no algorithm can tell from the noise rate whether a fixed circuit remembers
% Theorem 8.2 (conditional abstract form of Theorem 8.3).
\begin{theorem}[Conditional window certification]
\label{lean:conditional}
Let $u:\mathbb N\rightharpoonup\mathbb N$ be partial recursive. Suppose there
is a total computable compiler $\gamma$ from partial-recursive program codes
to inputs such that $u(\gamma(e))\downarrow$ if and only if $e(0)\downarrow$.
Let $c\in\mathbb Q$, $c>0$, let $L_0:\mathbb N\to\mathbb N$, and let
$I_n=(a_n,b_n)$ be pairwise disjoint open rational windows with
$0<a_n<b_n<1$. More precisely, for $n\ne m$, assume
$b_n\le a_m$ or $b_m\le a_n$.
Let $p_{n,L,t}\in\mathbb Q[X]$ for $n,L,t\in\mathbb N$.
Let $r:\mathbb N\to\mathbb Q$ be a surjection, and assume a total
computable Boolean test $A(n,L,t,k)$ satisfying exactly
\[
 A(n,L,t,k)=\mathrm{true}\quad\Longleftrightarrow\quad
 \bigl(L_0(n)\le L\ \land\ t\le cL\ \land\ a_n<r(k)<b_n
       \ \Longrightarrow\ p_{n,L,t}(r(k))\ge 7/8\bigr).
\]
Here all size and time variables are natural numbers; $t\le cL$ is
rational comparison and hence equivalent to $t\le\lfloor cL\rfloor$.
No separate computability assumptions on $L_0,a,b,p,r$ are needed:
the displayed exact test is the computational hypothesis.
Assume explicitly the following two physical premises of this paper:
\begin{enumerate}
\item (Nonhalting storage.) If $u(n)\uparrow$, then, for every $L,t\in\mathbb N$
and $q\in\mathbb Q$, whenever $L\ge L_0(n)$, $t\le cL$, and $q\in I_n$,
we have $p_{n,L,t}(q)\ge15/16$.
\item (Halting forgetting.) If $u(n)\downarrow$, there is $T\in\mathbb N$
such that for every $L,t\in\mathbb N$ and $q\in I_n\cap\mathbb Q$,
$t\ge T$ implies $p_{n,L,t}(q)\le9/16$.
\end{enumerate}
Define $P_{\mathbb Q}(n)$ to mean
\[
 \forall L,t\in\mathbb N\ \forall q\in\mathbb Q,\quad
 L\ge L_0(n),\ t\le cL,\ a_n<q<b_n
 \ \Longrightarrow\ p_{n,L,t}(q)\ge7/8.
\]
Define $P_{\mathbb R}(n)$ by replacing $q\in\mathbb Q$ by $x\in\mathbb R$,
coercing the endpoints and $7/8$ to $\mathbb R$, and evaluating the
coefficientwise image of $p_{n,L,t}$ in $\mathbb R[X]$ at $x$; keep
$t\le cL$ as the same rational comparison.
Then all four conclusions hold:
\[
 \begin{gathered}
 \forall n,\quad P_{\mathbb Q}(n)\ \Longleftrightarrow\ u(n)\uparrow;\\
 P_{\mathbb Q}\text{ is }\Pi^0_1\text{-complete};\\
 \forall n,\quad P_{\mathbb Q}(n)\ \Longleftrightarrow\ P_{\mathbb R}(n);\\
 P_{\mathbb Q}\text{ is not a computable predicate}.
 \end{gathered}
\]
For precision, $S\subseteq\mathbb N$ is $\Pi^0_1$ here when there exists
a total computable Boolean function $R:\mathbb N\times\mathbb N\to\{0,1\}$
with $S(n)\Longleftrightarrow\forall s\,R(n,s)=1$; completeness means that
$P_{\mathbb Q}$ is $\Pi^0_1$ and every such $S$ many-one reduces to it
by a total computable map $\mathbb N\to\mathbb N$.
\end{theorem}
\begin{proof}
The physical inputs are explicit hypotheses. Enumerating the exact tests over $(L,t,k)$ puts
$P_{\mathbb Q}$ in $\Pi^0_1$. Nonhalting storage gives every required
inequality. If the input halts, take the rational midpoint of its window,
the time $T$ from halting forgetting, and a size $L\ge L_0(n)$ with
$T\le cL$; this violates the certificate since $9/16<7/8$. Thus the
certified set is exactly the nonhalting set of $u$. The computable
compiler, together with the usual partial-recursive search for a failed
Boolean test, gives many-one $\Pi^0_1$-hardness and undecidability.
Finally, continuity of a real polynomial and density of the rationals in
every open real interval show that a strict real violation has a rational
violation; coefficient coercion preserves rational evaluation. This proves
the equivalence of rational and real certification.
\end{proof}

